\documentclass[aspectratio=169,11pt]{beamer}

% ── Packages ──────────────────────────────────────────────────────────────────
\usepackage[T1]{fontenc}
\usepackage[utf8]{inputenc}
\usepackage{lmodern}
\usepackage{microtype}
\usepackage{tikz}
\usetikzlibrary{shapes.rounded,arrows.meta,positioning,shadows,calc}
\usepackage{xcolor}

% ── Colour palette ────────────────────────────────────────────────────────────
\definecolor{Teal}      {HTML}{0D9488}
\definecolor{TealLight} {HTML}{CCFBF1}
\definecolor{Black}     {HTML}{111827}
\definecolor{Gray}      {HTML}{6B7280}
\definecolor{White}     {HTML}{FFFFFF}

% ── Beamer reset ──────────────────────────────────────────────────────────────
\usetheme{default}
\usecolortheme{default}
\setbeamertemplate{navigation symbols}{}
\setbeamertemplate{headline}{}
\setbeamertemplate{footline}{}
\setbeamercolor{background canvas}{bg=White}
\setbeamersize{text margin left=0.9cm, text margin right=0.7cm}

% ── Frametitle: large bold black (reference image style) ──────────────────────
\setbeamertemplate{frametitle}{%
  \vspace*{0.40cm}%
  \textbf{\LARGE\color{Black}\insertframetitle}%
  \par\vspace*{0.10cm}%
}

% ── Itemize bullets: teal filled circles ─────────────────────────────────────
\setbeamertemplate{itemize item}{%
  \tikz[baseline=-0.5ex]\filldraw[Teal,draw=Teal](0,0)circle(3.2pt);%
}
\setbeamertemplate{itemize subitem}{%
  \tikz[baseline=-0.5ex]\filldraw[Teal,draw=Teal](0,0)circle(2.2pt);%
}
\setbeamercolor{itemize item}{fg=Teal}
\setbeamercolor{itemize subitem}{fg=Teal}
\setlength{\leftmarginii}{1.3em}

% ── TikZ node styles ─────────────────────────────────────────────────────────
\tikzset{
  pipebox/.style={
    draw=Teal, fill=White,
    rounded corners=9pt,
    line width=1.4pt,
    minimum width=4.6cm, minimum height=0.70cm,
    font=\small\bfseries\color{Black},
    align=center,
    drop shadow={fill=black!10, shadow xshift=1.5pt, shadow yshift=-1.5pt}
  },
  fatarrow/.style={
    -{Stealth[length=10pt, width=9pt, inset=0pt]},
    line width=4pt,
    color=Teal,
    shorten >=2pt, shorten <=2pt
  },
  callout/.style={
    draw=Teal, fill=TealLight,
    rounded corners=6pt,
    line width=1.1pt,
    inner xsep=5pt, inner ysep=4pt,
    font=\scriptsize\bfseries\color{Black},
    align=center, text width=1.7cm
  }
}

% ── Goal box: text width accounts for inner sep (10pt each side = 20pt) ───────
\newcommand{\goalbox}[1]{%
  \noindent
  \begin{tikzpicture}
    \node[
      draw=Teal, fill=TealLight,
      rounded corners=9pt,
      line width=1.5pt,
      inner sep=9pt,
      text width=\dimexpr\linewidth-20pt\relax,
      font=\small\color{Black}
    ]{\textbf{Goal:}\ #1};
  \end{tikzpicture}%
}

% ══════════════════════════════════════════════════════════════════════════════
\begin{document}
% ══════════════════════════════════════════════════════════════════════════════


% ─────────────────────────────────────────────────────────────────────────────
%  SLIDE 1 — Medical Report Ingestion & OCR Pipeline
% ─────────────────────────────────────────────────────────────────────────────
\begin{frame}{Medical Report Ingestion \& OCR Pipeline}

  \begin{columns}[T]

    %% ── Left: bullets + goal box ─────────────────────────────────────────────
    \begin{column}{0.54\textwidth}

      {\small\itshape\color{Gray}%
        Convert raw scanned PDFs into raw OCR text stored in Supabase ---
        ready for downstream Gemini extraction.}

      \vspace{0.18cm}

      \begin{itemize}
        \setlength{\itemsep}{4pt}

        \item \textbf{Image preprocessing} on each PDF page
          \begin{itemize}
            \setlength{\itemsep}{1pt}
            \item Grayscale conversion and Gaussian denoising
            \item Adaptive thresholding and deskew via OpenCV
          \end{itemize}

        \item \textbf{Tesseract OCR} extracts raw text from images
          \begin{itemize}
            \setlength{\itemsep}{1pt}
            \item Page-by-page extraction with per-page confidence score
            \item Full raw text concatenated across all pages
          \end{itemize}

        \item \textbf{OCR text stored in Supabase}
          \begin{itemize}
            \setlength{\itemsep}{1pt}
            \item Stored in \texttt{medical\_reports.ocr\_text}
            \item Idempotent write --- safe to re-run on the same report
          \end{itemize}

      \end{itemize}

      \vspace{0.18cm}
      \goalbox{Convert scanned PDFs into raw OCR text stored in Supabase ---
        ready for Gemini LLM structured extraction.}

    \end{column}

    %% ── Right: pipeline diagram ───────────────────────────────────────────────
    \begin{column}{0.43\textwidth}
      \centering\vspace{0.10cm}

      \begin{tikzpicture}[node distance=0.46cm]

        \node[pipebox]                (pdf)  {PDF Upload};
        \node[pipebox, below=of pdf]  (pre)  {Image Preprocessing};
        \node[pipebox, below=of pre]  (ocr)  {Tesseract OCR};
        \node[pipebox, below=of ocr]  (norm) {Normalize Text};
        \node[pipebox, below=of norm] (db)   {Supabase Storage};

        \draw[fatarrow] (pdf)  -- (pre);
        \draw[fatarrow] (pre)  -- (ocr);
        \draw[fatarrow] (ocr)  -- (norm);
        \draw[fatarrow] (norm) -- (db);

        % Callout on OCR node
        \node[callout, anchor=west, xshift=0.24cm] (ann) at (ocr.east)
          {confidence\\score tracked};
        \draw[-{Stealth[length=5pt,width=4pt]}, Teal, line width=0.9pt]
          (ann.west) -- (ocr.east);

      \end{tikzpicture}

    \end{column}

  \end{columns}

\end{frame}


% ─────────────────────────────────────────────────────────────────────────────
%  SLIDE 2 — Gemini LLM Extraction Pipeline
% ─────────────────────────────────────────────────────────────────────────────
\begin{frame}{Gemini LLM Extraction Pipeline}

  \begin{columns}[T]

    %% ── Left ──────────────────────────────────────────────────────────────────
    \begin{column}{0.54\textwidth}

      {\small\itshape\color{Gray}%
        Gemini reads raw OCR text and extracts perfectly structured
        lab records --- handling noisy, varied report formats.}

      \vspace{0.18cm}

      \begin{itemize}
        \setlength{\itemsep}{4pt}

        \item \textbf{OCR text sent to Gemini} as a structured prompt
          \begin{itemize}
            \setlength{\itemsep}{1pt}
            \item System prompt encodes 10 clinical extraction rules
            \item User prompt wraps the full raw OCR text
          \end{itemize}

        \item \textbf{Structured JSON response} returned by Gemini
          \begin{itemize}
            \setlength{\itemsep}{1pt}
            \item Fields: \texttt{test\_name, value, unit, abnormal\_flag}
            \item Validated via Pydantic before any DB write
          \end{itemize}

        \item \textbf{Retry with exponential back-off} (no model switching)
          \begin{itemize}
            \setlength{\itemsep}{1pt}
            \item Up to 3 attempts on \texttt{gemini-3.1-pro-preview}
            \item Delay doubles each attempt: 2\,s $\to$ 4\,s $\to$ 8\,s
          \end{itemize}

        \item \textbf{Normalize \& idempotent upsert} to Supabase
          \begin{itemize}
            \setlength{\itemsep}{1pt}
            \item Canonical units, numeric cleaning, ref-range parsing
          \end{itemize}

      \end{itemize}

      \vspace{0.18cm}
      \goalbox{Gemini transforms ambiguous OCR text into structured,
        normalized lab records --- far beyond what any regex can handle.}

    \end{column}

    %% ── Right ─────────────────────────────────────────────────────────────────
    \begin{column}{0.43\textwidth}
      \centering\vspace{0.10cm}

      \begin{tikzpicture}[node distance=0.46cm]

        \node[pipebox]                   (src)    {OCR Text (Supabase)};
        \node[pipebox, below=of src]     (prompt) {Gemini Prompt Builder};
        \node[pipebox, below=of prompt]  (llm)    {Gemini 3.1-pro-preview};
        \node[pipebox, below=of llm]     (json)   {JSON Parsing (Pydantic)};
        \node[pipebox, below=of json]    (norm)   {Normalize \& Validate};
        \node[pipebox, below=of norm]    (db)     {Supabase Upsert};

        \draw[fatarrow] (src)    -- (prompt);
        \draw[fatarrow] (prompt) -- (llm);
        \draw[fatarrow] (llm)    -- (json);
        \draw[fatarrow] (json)   -- (norm);
        \draw[fatarrow] (norm)   -- (db);

        % Retry annotation — loop on LLM node
        \node[callout, anchor=west, xshift=0.24cm] (ann2) at (llm.east)
          {3 retries\\exp.\ back-off};
        \draw[-{Stealth[length=5pt,width=4pt]}, Teal, line width=0.9pt]
          (ann2.west) -- (llm.east);

      \end{tikzpicture}

    \end{column}

  \end{columns}

\end{frame}

\end{document}
